\Chapter{INTRODUCTION}\label{sec:Introduction}

Domain specific modeling languages for convex optimization such as CVXPY \cite{diamond2016cvxpy} and
YALMIP \cite{lofberg2004yalmip} allow users to formulate optimization problems in a natural
mathematical syntax. Through a process called \emph{canonicalization}, these high-level 
problem descriptions are converted to a \emph{standard conic form} used by many
general-purpose convex solvers
\BEQ \label{e:standard-form}
\begin{array}{ll}
\mbox{minimize} & c^T x + d\\
\mbox{subject to} & b-Ax \in K
\end{array}
\EEQ
where $x \in \reals^n$ is the decision variables,
and the \emph{problem data} are the matrix $A \in \reals^{m \times n}$, the vectors $b \in \reals^m$
and $c \in \reals^n$, the scalar $d \in \reals$ and the non-empty closed convex cone $K \subseteq \reals^m$.
Almost all of optimization practice can be divided into two groups, \emph{modeling} and 
\emph{solving}. The modeler's role is to formulate problems based on real-world business
or engineering applications.
The solver's role is to find an optimal solution given the data of a problem instance.
Modeling languages serve as the communication bridge between these two groups.
When improvements in solvers are made, modeling languages must adapt to allow users
to benefit from them. Conversely, it is through natural formulations that modelers can 
express their requirements and help motivate solver improvements.

While the time taken to solve an optimization problem is often dominated by the solver,
the time taken to \emph{canonicalize} (or \emph{compile}) a problem can also be significant,
especially for large problems.
In the case of CVXPY, the \emph{compilation time} largely lies in \emph{matrix stuffing},
the final step of canonicalization where symbolic expressions trees are converted to 
explicit numerical matrices.
Other modeling languages take different approaches to form the problem data, we review these
\emph{coefficient extraction} approaches in \S\ref{sec:coefficient-extraction}.

\paragraph{Atoms.}
Modeling languages for convex optimization are built on top of a library of
functions which can be used to describe problems.
These functions are known as \emph{atoms} in CVXPY since they are the fundamental
building blocks of the modeling language \cite{grant2004disciplined}.
Atoms have known attributes such as curvature (affine, convex, concave), monotonicity 
(increasing, decreasing, non-monotone), and sign (positive, negative, unknown), which are used
during canonicalization (we show an example in \S\ref{sec:canonicalization}).

\paragraph{Parameterized conic form.}
The data inputs in a problem instance can either be fixed or varying.
In CVXPY, varying data are referred to as \emph{parameters} and
fixed data are referred to as \emph{constants}.
Parameters are useful for specifying a \emph{family} of problems where the 
structure is fixed \cite{schaller2022cvxpygen}.
A \emph{parameterized cone program} extends the conic form in \eqref{eq:standard-form}
by letting the problem data depend on a parameter vector $\theta \in \reals^p$:
\BEQ \label{eq:param-cone}
\begin{array}{ll}
\text{minimize}   & c(\theta)^T x + d(\theta) \\
\text{subject to} & b(\theta) - A(\theta) x \in \mathcal K.
\end{array}
\EEQ
When the problem data $c, d, A, b$ are \emph{affine with respect to} $\theta$, we can
form the matrices $Q \in \reals^{(n+1) \times (p+1)}$ and
$\mathcal R \in \reals^{m \times (n+1) \times (p+1)}$ where
the problem data is obtained with
% convert this to keep the style consistent
\begin{equation} \label{eq:param-maps}
    \begin{bmatrix} c \\ d \end{bmatrix} = Q \cdot \tilde{\theta},
\qquad\text{and}\qquad
    [A \quad b] = \sum_{i=1}^{p+1} \mathcal{R}[:,:,i]\cdot \tilde{\theta}_i,
\end{equation}
where $d$ is the parameter-dependent constant offset of the objective
and $\tilde{\theta} = (\theta, 1) \in \reals^{p+1}$ (an extended parameter vector).

\paragraph{Parameter maps.}
During canonicalization, CVXPY will form the objective and constraints 
\emph{parameter maps} $Q$ and $\mathcal R$.
On subsequent solves, CVXPY only needs to update the parameter values
and apply the parameter maps in \eqref{param-maps}.
The benefits of the parameter maps is that it is completely separate from
the expression tree representation (which can be expensive when there are many nodes
in the problem formulation).
It can be viewed as a way to \emph{cache} parameter dependence in an optimization model.
% TODO make this a bit more polished, but the idea is there mostly.
We describe how the parameter map is computed in \S\ref{sec:matrix-stuffing}.

\paragraph{Expression trees.}
An expression is recursively defined as an atom evaluated at a \emph{subexpression}.
The subexpression can either be a variable, a parameter, or another expression.
Similarly to atoms, expressions can also be classified by curvature, monotonicity, and sign
which are determined by simple \emph{composition rules} \cite{grant2006disciplined}.
% define expressions and show an example of an expression tree.
An expression tree whose atoms are all affine is known as an \emph{affine expression tree}.

\paragraph{Matrix stuffing.}
% TODO we need to rewrite this paragraph now. I have some ideas but to revisit later.
% maybe talk about matrix stuffing with parameters here.
Although the original user specified problem formulation may contain nonlinearities,
after conversion to standard form, the \emph{canonicalized problem} will
only contain affine expressions.
The convex cone $\mathcal{K}$ (typically a cartesian product
of \emph{standard cones}) is already computed and known.
The remaining problem data are the \emph{constraint matrix} $A$
and the vectors $b$ and $c$.
The process to compute $b$ is relatively straightforward, for example,
one can \emph{evaluate} the affine expression tree at the origin.
The process to obtain $c$ is a special case of the process to obtain $A$.
For these reasons, we will often refer to matrix stuffing as the explicit task
of forming the $A$ matrix.
We review current CVXPY matrix stuffing approaches in \S\ref{sec:matrix-stuffing}.

\paragraph{Contributions.}
We make the following contributions:
\begin{enumerate}
  \item We propose an alternative method for matrix stuffing that doesn't rely on a parameter map.
  Instead, we treat matrix stuffing as a special case of sparse Jacobian evaluation which naturally
  separates the symbolic structure of the Jacobian from its numerical values.
  \item We use the affine atoms from the \texttt{SparseDiffengine} package and its Python bindings
  to implement a backend for CVXPY.
  \item We benchmark our method against existing approaches and make improvements to CVXPY's codebase
  by removing a deprecated backend.
\end{enumerate}
All code is open source and available at
\url{https://github.com/cvxpy/cvxpy}.

\paragraph{Thesis outline.}
The rest of this thesis begin with a review of the literature on modeling languages,
canonicalization, and matrix stuffing in \S\ref{sec:RevLitt}.
In \S\ref{sec:Theme1} we describe a set of relevant affine atoms and their implementation
in a sparse differentiation package.
% should we have another chapter for the CVXPY interface and problem class and handling parameters?
In \S\ref{sec:Theme2} we present benchmark results comparing our method
to existing approaches.
Finally, in \S\ref{sec:Conclusion} we summarize our contributions and
discuss future work.
